\section{Related Work}

\paragraph{Multi-Agent Systems (MAS) versus Single-Agent Systems (SAS)}

Understanding the difference between single-agent and multi-agent systems remains central to characterizing architectural effects. Following \citet{tran2025multiagent} and \citet{guo2024large}, we define a \textbf{Single-Agent System} as one that features a solitary reasoning locus: all perception, planning, and action occur within a single sequential loop controlled by one LLM instance, even when employing tool use \citep{yao2023react}, self-reflection \citep{shinn2023reflexion}, or chain-of-thought (CoT) reasoning \citep{wei2022emergent}. Critically, self-reflection mechanisms do not constitute multi-agent collaboration, as they operate within a single decision-making locus \citep{weng2023llmagent}. A \textbf{Multi-Agent System} comprises multiple LLM-backed agents communicating through structured message passing, shared memory, or orchestrated protocols \citep{xi2025rise}. MAS architectures vary by topology: \textit{Independent} systems aggregate isolated outputs; \textit{Decentralized} enable peer-to-peer exchange \citep{du2023improving}; \textit{Centralized} route through orchestrators \citep{hong2024metagpt}; \textit{Hybrid} combine hierarchical control with lateral communication \citep{dang2025evolving}. MAS evaluation has moved beyond early assumptions of uniform superiority \citep{li2024more, qian2024scaling} towards a more differentiated understanding driven by domain complexity. Recent surveys characterize collaboration mechanisms across coordination protocols \citep{tran2025multiagent} and agent profiling patterns \citep{guo2024large}. However, there exist empirical challenges: \citet{gao2025single} show benefits diminish as base models improve, with frontier models often outperforming teams; \citet{cemri2025multi} identify 14 failure modes (Cohen's Kappa=0.88); \citet{zhang2025maas} achieve comparable performance at 6-45\% cost through dynamic architecture search; and \citet{anthropic2024multiagent} report agents consume 15$\times$ more tokens. Theoretical foundations from \citet{sumers2024cognitive} propose cognitive architectures contextualizing agents within AI's broader history. The question of \textit{when} multi-agent coordination provides value over single strong models with tool use remains empirically open, with \citet{qian2024scaling}'s proposed scaling laws showing no significant universal pattern \citep{wang2024survey}, motivating our systematic evaluation.

\paragraph{Agentic Tasks and Benchmarks}

We define \textit{agentic tasks} following \citet{zhu2025establishing} as requiring: (1) sustained multi-step environment interactions, (2) iterative information gathering under partial observability, and (3) adaptive strategy refinement from feedback, differentiating tasks like web browsing \citep{wei2025browsecomp, zhou2024webarena}, financial trading \citep{bigeard2025finance}, software engineering \citep{jimenez2023swebench}, and planning \citep{dagan2024plancraft} from static benchmarks. \textit{Non-agentic tasks} evaluate single-shot inference without environmental interaction: GSM8K \citep{cobbe2021gsm8k} (direct chain-of-thought math), MMLU \citep{hendrycks2020mmlu} (parametric knowledge), HumanEval \citep{chen2021humaneval} (specification-complete coding), and SQuAD \citep{rajpurkar2016squad} (single-pass comprehension). On non-agentic benchmarks, multi-agent systems show monotonic improvement through ensemble effects (89\% on HumanEval with five agents), as voting corrects errors without sequential compounding \citep{kapoor2024agents}. This distinction matters: in agentic settings, coordination overhead scales with interaction depth, agents operate on divergent world states (34\% overlap after 10 interactions), and errors cascade rather than cancel \citep{kapoor2024agents}. \citet{zhu2025establishing} introduce the Agentic Benchmark Checklist addressing flaws causing 100\% relative misestimation. Evolution spans \citet{liu2023agentbench}'s 8-environment evaluation (4k-13k responses) to specialized frameworks: \citet{jimenez2023swebench} (GitHub resolution), \citet{zhou2024webarena} (812 web tasks), \citet{xu2024theagentcompany} (30\% autonomous completion), and \citet{paglieri2024balrog} (vision-based RL). \citet{yao2023react} formalizes reasoning-acting synergy; \citet{weng2023llmagent} characterizes agents requiring planning, memory, and tools; \citet{kapoor2024agents} reveals narrow accuracy focus without cost metrics yields needlessly complex agents. \textcolor{black}{We note that established agentic benchmarks such as SWE-bench \cite{jimenez2023swe}, WebArena, and Tau-bench already embody these evaluation properties, as discussed in recent survey work \cite{yehudai2025survey}. Our contribution is not the formalization of these properties per se, but rather their systematic application as experimental controls across five coordination architectures and nine models from three LLM families, enabling the first quantitative characterization of how coordination benefit scales with model capability.} Tasks showing MAS advantages in single-shot settings often exhibit opposite patterns under genuine interaction, indicating architectural benefits are task-contingent, motivating our isolation of coordination effects across diverse agentic domains.

\paragraph{Scaling Laws and Coordination Mechanisms}

Understanding performance scaling in multi-agent systems requires distinguishing collaborative scaling from neural scaling laws. While neural scaling follows power laws requiring million-fold parameter increases for significant trends \citep{kaplan2020scaling}, collaborative scaling exhibits logistic growth patterns emerging at substantially smaller scales \citep{qian2024scaling}. \citet{chen2024compound} explore whether increased LLM calls alone drive performance, finding compound inference systems follow distinct scaling behaviors from single-model training. However, \citet{wang2024survey} note collaborative scaling shows no significant universal pattern, suggesting domain-specific rather than general laws. Coordination mechanisms critically determine whether collaboration amplifies or degrades performance: \citet{hong2024metagpt} introduce meta-programming workflows mitigating hallucination cascades; \citet{chen2023agentverse} demonstrate emergent behaviors through structured interactions; \citet{wu2024autogen} provide general multi-agent frameworks. Recent work reveals architecture-task alignment matters more than team size: \citet{zhang2025maas} achieve superior performance at 6-45\% cost through query-dependent configurations; \citet{dang2025evolving} show puppeteer orchestration improvements stem from compact cyclic structures; \citet{du2023improving} demonstrate peer-to-peer debate effectiveness depends on task decomposability, with \citet{smit2023should} further showing that multi-agent debate does not reliably outperform single-agent strategies such as self-consistency, suggesting benefits are highly task- and hyperparameter-sensitive. These findings collectively indicate coordination benefits arise from matching communication topology to task structure not from scaling the number of agents, establishing the foundation for principled architectural design rather than heuristic ``more agents is better'' approaches.